%==============================================================================
\chapter{Introdução}\label{introducao}
%==============================================================================

A tecnologia digital tornou-se pervasiva em domínios diversificados, inclusive no educacional. Neste domínio, estudantes, professores e gestores utilizam tecnologias digitais para desempenharem uma série de tarefas: ler, escrever, calcular, planejar, avaliar etc. Quando se trata da gestão escolar, é comum a adoção de ferramentas de escritório. Além disso, há relatos de sistemas de software que apoiam a contratação e a alocação de professores em escolas, o acompanhamento da vida escolar do estudante pela família, a organização das atividades escolares, o acompanhamento de unidades escolares, assim como a gestão de dados de planejamento escolar, registro escolar e estrutura de funcionamento \cite{OliveiraMelo2020}


Ao investigarem o domínio educacional, com ênfase para a gestão escolar, Oliveira e Melo (2020) propuseram o aplicativo móvel ParecerEdu para apoiar a emissão de pareceres nos anos iniciais do ensino fundamental. Segundo os autores, ps trabalhos relacionados apresentaram o uso de tecnologias digitais no apoio à gestão escolar. Contudo, nenhum especificava seu processo de Engenharia de Software nem como ocorreu a participação dos usuários em sua concepção. Ademais, nenhuma dessas soluções é voltada especificamente à elaboração de pareceres descritivos pelos professores, a partir de dados de acompanhamento dos estudantes no dia a dia da sala de aula \cite{OliveiraMelo2020}.	
	
Nesse contexto, foi criado o ParecerEdu. Este tem como propósito auxiliar professores dos anos iniciais do ensino fundamental na geração de pareceres a partir de dados de acompanhamento de estudantes na sala de aula.  Para fazê-lo, o aplicativo possibilita aos professores, para cada uma de suas turmas, definir critérios de avaliação em diferentes áreas do conhecimento e avaliar determinada turma a partir desses critérios, indicando, para cada aluno, se um critério foi atingido, atingido parcialmente ou não foi atingido. Gráficos são gerados para apresentar o desempenho da turma e também o desempenho individual dos alunos.

Uma primeira versão do aplicativo foi compartilhada com professoras dos anos iniciais do ensino fundamental das escolas estaduais Arthur Hormain e Oswaldo Dornelles de Alegrete/RS. Contudo, seu uso precisa ser acompanhado, o que deve gerar novos requisitos, levando à evolução desse sistema. Além disso, Otokovieski, Melo e Crespo (2020) realizaram uma avaliação de usabilidade do aplicativo, identificando problemas com alto grau de severidade que devem ser corrigidos.

Segundo Sommervile (2018), um software necessita evoluir para se manter útil e, assim, as mudanças no sistema devem ser feitas para garantir que ele continue a agregar valor. A evolução em um contexto geral acontece por inúmeros motivos, como mudanças nas empresas e nas expectativas do usuário. Neste trabalho, a principal motivação é que a experiência dos usuários com o software tendem a gerar novos requisitos (WANG; LIU; WANG, 2012; CARNEIRO et al., 2020). Além disso,  as adaptações a plataformas devem ser feitas para que o software seja adequado a tecnologias atuais.  
Nesse contexto, numa perspectiva de Engenharia de Software, pergunta-se: Como a experiência do usuário pode influenciar na evolução do app ParecerEdu? Sua arquitetura comporta o desenvolvimento de novos requisitos? De que modo evoluir o app para que ele possa ser utilizado em diferentes plataformas? 
O objetivo geral deste Trabalho de Conclusão de Curso é evoluir o aplicativo app ParecerEdu a partir da experiência do usuário com esse software. Como objetivos específicos, tem-se:

\begin{itemize}
  \item Investigar o domínio educacional;
\item Evoluir o aplicativo ParecerEdu;
\item Integrar práticas de Interação Humano-computador à Engenharia de Software;
\item  o desenvolvimento de front-end de sistemas móveis.

\end{itemize}

Como referências para embasar este trabalho, foi investigado nas bases ACM Digital Library, IEEE Xplore e Scopus, afim de encontrar trabalhos acadêmicos que levam em conta a evolução do software observando experiência do usuário e o que usa-se de apoio para realizar essa evolução, esta etapa é descrita no capítulo X. Logo é apresentada a metodologia realizada no capítulo X, abordando o processo realizado para elaboração deste estudo, assim como ferramentas usadas e resultados obtidos a partir disso.
Como referências para embasar este trabalho, foi investigado nas bases ACM Digital Library, IEEE Xplore e Scopus, afim de encontrar trabalhos acadêmicos que levam em conta a evolução do software observando experiência do usuário e o que usa-se de apoio para realizar essa evolução, esta etapa é descrita no capítulo X. Logo é apresentada a metodologia realizada no capítulo X, abordando o processo realizado para elaboração deste estudo, assim como ferramentas usadas e resultados obtidos a partir disso.



